%% ===========================================================================
\FloatBarrier
\section{Conclusions}
\label{sec:conclusions}
%% ===========================================================================

The formulation combines piezoelectric and flexoelectric coupling with
fractional memory in a $C^1$ Bogner--Fox--Schmit discretization. The prescribed window second moment parametrizes the assumed local
gradient stiffness, while modal L1
integration retains the reduced mechanical stiffness history. The principal findings for the prescribed parameters are:

\begin{itemize}
\item The window second-moment tensor $M$ defines the internal length
  parameter in the adopted local gradient-energy model. Windows with the same
  $M$ give identical solutions within this second-moment approximation; this
  does not establish equivalence of their original spatial-averaging kernels,
  particularly where a nominal window extends outside the strip. The local-model numerical design map spans coupling from
$3.6\times10^{-6}$ to $1.0\times10^{-2}$ across all 64 cases, while
the 25 cases satisfying the necessary mid-plane window-fit condition
$2a_2\le h$ span $1.28\times10^{-4}$ to $1.01\times10^{-2}$.
Neither subset, without a boundary-truncated averaging kernel or
experiment, establishes a calibrated finite-specimen design law. Rotating an anisotropic window changes the extensional
  coupling by $14$--$17\,\%$ between the aligned and transverse positions,
  because the gradient stiffness samples different directional strain gradients.
\item For a prescribed bending curvature, the strain gradient is set by the
  curvature and the electrode/geometry response; the two coupling mechanisms exhibit different thickness dependences: the
  flexoelectric bending coupling grows by a factor $24$ between
  \SI{10}{\micro\meter} and \SI{1}{\micro\meter}, while the piezoelectric
  extensional coupling stays constant to $0.07\,\%$ and its quasi-static value,
  $0.5727$, lies within $0.5\,\%$ of the closed form $0.5752$.
\item The assumed flexoelectric coefficient $\mu$ influences the magnitude of
  the bending coupling: at
  the assumed $\mu=10^{-6}$\,\si{\coulomb\per\meter} the ratio of the two
  couplings spans $2.9\times10^{-4}$ to $7.0\times10^{-3}$ and reaches
  $0.1\,\%$ only below $h=5.3\,\si{\micro\meter}$.
\item The Caputo order $\alpha$, which specifies the fractional mechanical memory response, controls the shape of the decay: the first-period
  equivalent damping index of the tip displacement rises from $0.065$ at
  $\alpha=1$ to a maximum of $0.105$ at $\alpha=0.6$ and then decreases
  to $0.085$ at $\alpha=0.4$ for
  the tested time scale. The fractional decay is nonexponential, and replacing
  the $\alpha=0.6$ early-decay measure with its integer-order value would change that index substantially: the fractional value is
  $62\,\%$ larger than the integer-order value.
\item The time parameter $\tau$ controls loss through the
  dimensionless frequency $\omega\tau$: at $\alpha=1$ and
  $\omega\tau=0.3$ the damping ratio is $0.15003$ on the condensed matrix
  against $0.0646$ on the mechanical-only one, a ratio of $0.431$, and the
  Kelvin--Voigt limit is reproduced to $0.02\,\%$.
\item Mesh refinement of the $C^1$ Bogner--Fox--Schmit method converges
  in the tested benchmarks; the gradient-free bar sequence has
  an observed overall order of $1.68$, increasing to $2.01$ between its
  finest two meshes. In the extensional coupling study, refinement from
  $48\times4$ to $96\times4$ elements changes the coupling factor by
  $5.9\times10^{-4}$ relative.
\end{itemize}

Within the adopted model, electrode choice ($0.4781$ against $0.1036$
at equal tip motion), thickness and gradient-length parameters govern the
reported sensing response. The design map includes configurations in which a
literal window would cross the specimen faces; those results quantify local
model sensitivity rather than calibrated physical device performance. Decay
interpretation depends jointly on $\alpha$, $\tau$ and the chosen memory
matrix. Quantitative application to a particular material requires calibration of
the assumed inputs and assessment of the two-dimensional approximation:
plane-stress coupling differs from the three-dimensional value by $0.3\,\%$.
Calibration of the window parameters and a three-dimensional extension are
needed for more general predictions.

\FloatBarrier
